%======================================================================
\section{Concluding Remarks}\label{sec:conclusion}

We have shown that a polarized extension-variable translation
preserves \kl{minimality} with respect
to the \kl{substitution preorder} in every \kl{logic} that contains the axiom $\axI$
and the translated formula. As a consequence, axiomatizations of
implicational logics over $\BCI$ by $\CL$-minimal formulas can
always be taken to consist of formulas of \kl{order} at most $4$, and
\cite[Open Problem~2]{NM21} has an affirmative answer. We conclude with
some questions that remain open.
\begin{enumerate}
\item Which intermediate logics are axiomatizable by $\CL$-minimal formulas
  \cite[Open Problem~1]{NM21}? By Corollary~\ref{cor:normal-form} and
  Proposition~\ref{prop:order3}, it suffices to consider formulas of order $4$.
\item Is there a minimality-preserving translation over $\FLe$
  with additive connectives, or over the non-commutative calculus
  $\FL$ (Section~\ref{sec:additive})?
\item Does the translation, possibly with the formulas $\dprem{u}$ listed in a different order, preserve
  \kl{minimality} over ticket entailment $\Tto$
  (Appendix~\ref{app:weak})?
\item What is the least number of \kl{antecedents} of order $3$ in a
  $\CL$-minimal axiom of order $4$ of a proper intermediate logic? By
  Proposition~\ref{prop:one-premise} it is at least $2$, whereas $\St(G')$
  has $11$ such \kl{antecedents}.
\end{enumerate}

\paragraph{Acknowledgements.} \plot{JSPS KAKENHI etc.}
